% GENERATED FILE. Edit canonical lessons, then run npm run generate:books.
% canonical-source-hash: fnv1a64:c40c1902496ded77
% canonical-lessons: LA-C48-nescio, LA-C48-male, LA-C48-malus, LA-C48-sed, LA-C48-aut

\chapter{Nescio, Male, Malus, Sed, Aut}
\label{ch:la-negatio}
\hlchaptermodality{48}

\begin{chapteropening}
\textbf{By the end of this chapter:} \emph{I can say that I do not know, that something went badly, and I can set one thing against another or offer a choice between two.}

The fifth word of the group, taken apart against the four before it.
\end{chapteropening}

\section[nesciō, nescīre]{nesciō — and the rule the whole chapter turns on: nōn stands in front of the verb}

\label{lesson:LA-C48-nescio}



\begin{quote}
Six recalls before the new word, at growing distances. They are not decoration: each one is the retrieval window this book measures, and each names a word far enough back that saying it is work.

\begin{itemize}
\raggedright
  \item \emph{Recall:} say \emph{piscis} — \textbf{R1}, one lesson back, before it decays
  \item \emph{Recall:} say \emph{sāl} — \textbf{R2}, the first real retrieval, five to fifteen lessons back
  \item \emph{Recall:} say \emph{bonum māne} — \textbf{R3}, durable, twenty to sixty lessons back
  \item \emph{Recall:} say \emph{Hesperia} — \textbf{R3}, durable again, a second word from the same distance
  \item \emph{Recall:} say \emph{caput} — \textbf{R4}, recognition at distance, eighty lessons back or more
  \item \emph{Recall:} say \emph{manus} — \textbf{R4}, recognition at distance, a second word from that far back
\end{itemize}
\end{quote}

\begin{cousinweb}
\textbf{nesciō} — \textquotedblleft{}\textbf{I do not know}.\textquotedblright{}

Look at the shape. It is \textbf{nē-} (\textquotedblleft{}not\textquotedblright{}) welded onto \textbf{sciō}, the verb the eight-verb chapter gave you. Latin fused the negative into the verb once and kept it: \emph{nesciō} is one word, and it is the word a Roman actually said.

\textbf{Now the rule that unlocks everything else.} For every other verb, Latin puts \textbf{nōn} — the word your very first lesson taught you as \textquotedblleft{}no\textquotedblright{} — \textbf{immediately in front of the finite verb}:

\begin{itemize}
\raggedright
  \item \textbf{nōn valeō} — \textquotedblleft{}I am not well.\textquotedblright{}
  \item \textbf{nōn intellegō} — \textquotedblleft{}I do not understand.\textquotedblright{}
  \item \textbf{nōn videō} — \textquotedblleft{}I do not see.\textquotedblright{}
\end{itemize}

That is the whole of it. No auxiliary, no change to the verb, no word order to learn. You have had \textbf{nōn} since your first day and every one of these verbs since the eight-verb chapter, and until now nothing told you that you could put them together.

English \textbf{nescience} is \emph{nesciō} borrowed whole. \textbf{Nice} is the same word by a long road: Latin \emph{nescius} (\textquotedblleft{}ignorant\textquotedblright{}) became Old French \emph{nice} (\textquotedblleft{}silly\textquotedblright{}), and seven centuries of drift turned an insult into a compliment.
\end{cousinweb}

\begin{grammarlens}[title={Grammar Lens: non goes in front of the verb}]
State it once, plainly, because everything else in this tranche leans on it.

\begin{quote}
\textbf{nōn} + a finite verb = the negative of that verb.
\end{quote}

\begin{itemize}
\raggedright
  \item \textbf{nōn valeō} — I am not well.
  \item \textbf{nōn intellegō} — I do not understand.
  \item \textbf{nōn sciō} — I do not know (and \textbf{nesciō} is the word a Roman preferred).
  \item \textbf{nōn possum} — I cannot. (The next chapter's verb, negated with this rule.)
\end{itemize}

Nothing else changes. Latin does not add an auxiliary, does not move the verb and does not have a second negative word to pair it with. One word, in front.
\end{grammarlens}

\subsection*{Your turn}
Say these aloud:

\begin{itemize}
\raggedright
  \item \emph{nesciō, nescīre}
  \item every word this chapter has given you so far, in order
\end{itemize}

\begin{tcolorbox}[breakable,colback=teal!4,colframe=teal!35!black,title={Before you move on}]
How does Latin negate an ordinary verb? (\textbf{nōn} immediately before it.) Say \textquotedblleft{}I do not understand.\textquotedblright{} (\emph{\textbf{Nōn intellegō}}.) And what two words are inside \emph{nesciō}? (\emph{\textbf{Nē-}} and \emph{\textbf{sciō}}.)
\end{tcolorbox}
\section[male]{male — the other end of the scale bene has had to itself all along}

\label{lesson:LA-C48-male}



\begin{quote}
Six recalls before the new word, at growing distances. They are not decoration: each one is the retrieval window this book measures, and each names a word far enough back that saying it is work.

\begin{itemize}
\raggedright
  \item \emph{Recall:} say \emph{nesciō} — \textbf{R1}, one lesson back, before it decays
  \item \emph{Recall:} say \emph{lūna} — \textbf{R2}, the first real retrieval, five to fifteen lessons back
  \item \emph{Recall:} say \emph{bonum vesperum} — \textbf{R3}, durable, twenty to sixty lessons back
  \item \emph{Recall:} say \emph{merīdiēs} — \textbf{R3}, durable again, a second word from the same distance
  \item \emph{Recall:} say \emph{manus} — \textbf{R4}, recognition at distance, eighty lessons back or more
  \item \emph{Recall:} say \emph{vēr} — \textbf{R4}, recognition at distance, a second word from that far back
\end{itemize}
\end{quote}

\begin{cousinweb}
\textbf{male} — \textquotedblleft{}\textbf{badly}.\textquotedblright{}

The \emph{bene} lesson gave you \textbf{bene}, \textquotedblleft{}well,\textquotedblright{} and left you with half a scale. This is the other half, and it is built the same way: an adverb ending in \textbf{-e}, formed from an adjective.

\begin{quote}
\textbf{Valeō.} — I fare well. \textbf{Male valeō.} — I fare badly. (Or, more usually: \textbf{nōn valeō}.)
\end{quote}

English keeps \emph{male-} as a prefix that means \textquotedblleft{}wrongly\textquotedblright{}: \textbf{mal}ice, \textbf{mal}ady, \textbf{mal}ign, \textbf{mal}practice, \textbf{mal}function. \textbf{Malaria} is Italian \emph{mal'aria}, \textquotedblleft{}bad air\textquotedblright{} — what people thought caused it before the mosquito was known. \textbf{Dismal} is \emph{diēs malī}, \textquotedblleft{}bad days,\textquotedblright{} the unlucky days of the medieval calendar.
\end{cousinweb}

\subsection*{Your turn}
Say these aloud:

\begin{itemize}
\raggedright
  \item \emph{male}
  \item every word this chapter has given you so far, in order
\end{itemize}

\begin{tcolorbox}[breakable,colback=teal!4,colframe=teal!35!black,title={Before you move on}]
What is the opposite of \emph{bene}? (\emph{\textbf{Male}}.) What two Latin words are inside \emph{dismal}? (\emph{\textbf{Diēs malī}} — bad days.)
\end{tcolorbox}
\section[malus, mala, malum]{malus — the adjective under male, and the word Rome punned on for the apple}

\label{lesson:LA-C48-malus}



\begin{quote}
Six recalls before the new word, at growing distances. They are not decoration: each one is the retrieval window this book measures, and each names a word far enough back that saying it is work.

\begin{itemize}
\raggedright
  \item \emph{Recall:} say \emph{male} — \textbf{R1}, one lesson back, before it decays
  \item \emph{Recall:} say \emph{stēlla} — \textbf{R2}, the first real retrieval, five to fifteen lessons back
  \item \emph{Recall:} say \emph{merīdiēs} — \textbf{R3}, durable, twenty to sixty lessons back
  \item \emph{Recall:} say \emph{habeō} — \textbf{R3}, durable again, a second word from the same distance
  \item \emph{Recall:} say \emph{vēr} — \textbf{R4}, recognition at distance, eighty lessons back or more
  \item \emph{Recall:} say \emph{aqua} — \textbf{R4}, recognition at distance, a second word from that far back
\end{itemize}
\end{quote}

\begin{cousinweb}
\textbf{malus} — \textquotedblleft{}\textbf{bad}.\textquotedblright{}

The adjective that \textbf{male} was cut from. It changes its ending to match its noun, exactly as \textbf{bonus} does in \emph{bonam noctem}: \emph{malus} (m.), \emph{mala} (f.), \emph{malum} (n.).

Now the coincidence that changed a story. Latin also had \textbf{mālum}, \textquotedblleft{}\textbf{apple},\textquotedblright{} borrowed from Greek \emph{mēlon} — a different word, with a long \emph{ā}, and in writing without macrons the two look identical. Genesis names no fruit at all: it says only \textquotedblleft{}the fruit of the tree.\textquotedblright{} When Latin translators wrote of the tree of \emph{mālum}, the pun with \emph{malum}, \textquotedblleft{}evil,\textquotedblright{} was there for the taking, and medieval art took it.

\textbf{Treat that as the standard account rather than as settled fact}: the pun is real and the timing is right, and no single document records anybody deciding it. English \textbf{malevolent}, \textbf{malefactor} and \textbf{malcontent} are all \emph{malus}.
\end{cousinweb}

\subsection*{Your turn}
Say these aloud:

\begin{itemize}
\raggedright
  \item \emph{malus, mala, malum}
  \item every word this chapter has given you so far, in order
\end{itemize}

\begin{tcolorbox}[breakable,colback=teal!4,colframe=teal!35!black,title={Before you move on}]
What does \emph{malus} mean, and what does \emph{mālum} mean? (\textbf{Bad}; \textbf{apple}.) What separates them in writing? (\textbf{The macron} — vowel length.)
\end{tcolorbox}
\section[sed]{sed — the word that sets one thing against another}

\label{lesson:LA-C48-sed}



\begin{quote}
Six recalls before the new word, at growing distances. They are not decoration: each one is the retrieval window this book measures, and each names a word far enough back that saying it is work.

\begin{itemize}
\raggedright
  \item \emph{Recall:} say \emph{malus} — \textbf{R1}, one lesson back, before it decays
  \item \emph{Recall:} say \emph{umbra} — \textbf{R2}, the first real retrieval, five to fifteen lessons back
  \item \emph{Recall:} say \emph{eō} — \textbf{R3}, durable, twenty to sixty lessons back
  \item \emph{Recall:} say \emph{vesper — vespere — sērus} — \textbf{R3}, durable again, a second word from the same distance
  \item \emph{Recall:} say \emph{aqua} — \textbf{R4}, recognition at distance, eighty lessons back or more
  \item \emph{Recall:} say \emph{pānis} — \textbf{R4}, recognition at distance, a second word from that far back
\end{itemize}
\end{quote}

\begin{cousinweb}
\textbf{sed} — \textquotedblleft{}\textbf{but}.\textquotedblright{}

You have met it once already without being told. The \emph{propediem tē vidēbō} lesson quoted Cicero's own letter: \emph{\textquotedblleft{}\textbf{Sed}, ut spērō, propediem tē vidēbō\textquotedblright{}} — \textquotedblleft{}\textbf{But}, as I hope, I shall see you before long.\textquotedblright{} That was a real sentence by a real Roman, and this is the word that opened it.

\textbf{Sed} joins two statements and marks the second as a correction or a contrast:

\begin{quote}
\textbf{Nōn valeō, sed nōn male valeō.} — I am not well, but I am not badly off.
\end{quote}

Here is the strange part. \emph{Sed} is one of the commonest words in surviving Latin prose and \textbf{not one Romance language kept it}. Spanish and Italian took \emph{magis} (\textquotedblleft{}more\textquotedblright{}) for \emph{mas} and \emph{ma}; French took \emph{mais} from the same \emph{magis}; Portuguese \emph{mas} likewise. A word Cicero could not write a page without vanished completely.
\end{cousinweb}

\subsection*{Your turn}
Say these aloud:

\begin{itemize}
\raggedright
  \item \emph{sed}
  \item every word this chapter has given you so far, in order
\end{itemize}

\begin{tcolorbox}[breakable,colback=teal!4,colframe=teal!35!black,title={Before you move on}]
What does \emph{sed} mean? (\textbf{But}.) Which Roman did you first meet it from? (\textbf{Cicero}, in the \emph{propediem} lesson.) Did any daughter language keep it? (\textbf{No} — they all took \emph{magis} instead.)
\end{tcolorbox}
\section[aut]{aut — and the second or, which asks a different question}

\label{lesson:LA-C48-aut}



\begin{quote}
Six recalls before the new word, at growing distances. They are not decoration: each one is the retrieval window this book measures, and each names a word far enough back that saying it is work.

\begin{itemize}
\raggedright
  \item \emph{Recall:} say \emph{sed} — \textbf{R1}, one lesson back, before it decays
  \item \emph{Recall:} say \emph{lūx} — \textbf{R2}, the first real retrieval, five to fifteen lessons back
  \item \emph{Recall:} say \emph{bonum vesperum} — \textbf{R3}, durable, twenty to sixty lessons back
  \item \emph{Recall:} say \emph{merīdiēs} — \textbf{R3}, durable again, a second word from the same distance
  \item \emph{Recall:} say \emph{pānis} — \textbf{R4}, recognition at distance, eighty lessons back or more
  \item \emph{Recall:} say \emph{Iānuārius} — \textbf{R4}, recognition at distance, a second word from that far back
\end{itemize}
\end{quote}

\begin{cousinweb}
\textbf{aut} — \textquotedblleft{}\textbf{or}.\textquotedblright{}

The word that lets you offer a choice, which until now you could not do:

\begin{quote}
\textbf{Aquam aut vīnum?} — Water or wine?
\end{quote}

Latin has \textbf{two} words here and they are not interchangeable. \textbf{Aut} offers an \textbf{exclusive} choice: one or the other, not both. \textbf{Vel} offers an \textbf{indifferent} one: either will do. \emph{Vel} is literally the imperative of \emph{volō} — \textquotedblleft{}choose, if you like\textquotedblright{} — which is the next chapter's verb.

Logicians still use the distinction. The inclusive-or of Boolean algebra is called \textbf{vel} after the Latin, and it is inclusive precisely because \emph{vel} was; exclusive-or is \emph{aut ... aut}. A grammatical nicety of a dead language is sitting inside every computer you have ever used.
\end{cousinweb}

\subsection*{Your turn}
Say these aloud:

\begin{itemize}
\raggedright
  \item \emph{aut}
  \item every word this chapter has given you so far, in order
\end{itemize}

\begin{tcolorbox}[breakable,colback=teal!4,colframe=teal!35!black,title={Before you move on}]
Which \emph{or} excludes, and which does not? (\emph{\textbf{Aut}} excludes; \emph{\textbf{vel}} does not.) What is \emph{vel} literally? (\textbf{The imperative of \emph{volō}} — choose if you like.)
\end{tcolorbox}
